\section{Related work}
\label{sec:related:work}


\paragraph*{Planar visibility task planning}
Much of the existing literature on \vistamp focuses on planar environments, where efficient and, in some cases, optimal solutions often exist. Planar \vistamp algorithms benefit from the simplicity and abundance of visibility algorithms \cite{g-vap-07}. Optimal algorithms for achieving and maintaining visibility of a single moving target with known or partially-predictable movements were introduced in \cite{lgbl-msmvmt-97}. Unpredictable moving targets were later addressed in \cite{MCTH-sbmpamvut-05}. Many more planar \vistamp problems were tackled, such as landmark-based navigation \cite{bmh-oplbnddvfvc-07}, simultaneous tracking and localization \cite{wlzhlf-ovmpfttl-14}, pursuit-evasion games \cite{bh-enetppegvc-08}, and multi-agent tracking \cite{rl-mrtdtts-16}.

\paragraph*{Spatial visibility task planning} 
Spatial \vistamp requires solving spatial visibility problems, some of which are computationally expensive and susceptible to floating point errors \cite{d-3dvasa-99}. \new{To plan trajectories for multi-target tracking, a Visibility Roadmap (\vir) based on \vi was established in \cite{lk-fmtvw-16, zhvml-c3dfrvuvi-19}. \vir assumes a 3D point robot with an omni-directional sensor, answering queries by sampling within high-\vi clusters, derived from either a grid or randomly sampled points. Conversely, our method supports complex systems by using a hierarchical decomposition to compute visibility relationships between different regions.} 
% In \cite{lk-fmtvw-16, zhvml-c3dfrvuvi-19} the visibility integrity (\vi) measure and a visibility roadmap (\vir) based on it is introduced and used to plan trajectories for following multiple targets.} 
\new{Another} \prm variant was used to compute efficient paths for exhaustively searching an object in a 3D environment \cite{smh-sbccrfo3de-05}. %\new{Both} algorithms assume and heavily rely on an omni-directional sensor model. 
%\new{The \vir papers also assume a point robot}. 
Similarly, a \prm graph was augmented to produce motion plans maintaining a target in the robot's \fov \cite{blcl-ppivuprm-10}. A different work uses \prm for surface disinfection using UV light \cite{mrph-ocpuvsd-21}. It constructs a roadmap while keeping track of the regions illuminated by the UV device in the sampled configurations. The roadmap is then used to plan a path covering the set of target surfaces. We compare our proposed algorithms to baseline \prm and \rrt planners based on this work. \new{A modified $\text{RRT}^*$ framework was deployed in \cite{byam-lprrtdsrccf-22} using various local planners to optimize motion paths across distinct cost functions. However, their visibility application is strictly limited to maintaining a target within the FOV, bypassing the problem of active target exploration.} Another work plans paths that maintain visibility with a target or avoid entering a sentry's visibility region \cite{pld-prpmpsvr-24}. 




%Unlike the planar case, many \vistamp problems previously mentioned, including \tvmp, do not yet have algorithms reliably capable of solving them.  


\paragraph*{Other spatial visibility tasks}
A large body of work is dedicated to local visibility tasks. Visual servoing \cite{ch-vsvt-08} is a widely-used technique for robot motion planning and control using a camera. Visual servoing problems include the eye-in-hand model in which a camera is attached to the robot and motion planning is performed where the goal state is a configuration in which the camera's view matches an input image. In the visual tracking problem \cite{bah-mp3dttao-07,ghnd-kmmtfvmc-11,lwlcz-udmavtcrehfcc-15,daop-mvdoceummvfi-23}, an adjustable camera or a camera-wielding robot must keep a moving target or set of targets within its \fov. While these problems share similarities with the \tvmp, they typically assume an initial state where the target is already within the \fov. Consequently, they focus on local reactive strategies to maintain visibility, making control decisions based on the target’s relative displacement and immediate workspace constraints. Algorithms developed for such tasks are thus unsuitable for global visibility tasks, and cannot be used as benchmarks in our experiments. Another set of \vistamp algorithms are designed only for controlled conditions such as robots in unconstrained environments or with a predefined set of allowable paths \cite{w-vbopmp-17}, e.g., on rails.

\paragraph*{Sampling-based motion planning}
The proposed algorithms, \visrrt and \visprm, follow the \sbmp paradigm in which a geometric graph $G=(V,E)$ is constructed in the robot's \cspace ~using random sampling. $V$ and $E$ are associated with robot configurations and continuous motions respectively, and an $(s,t)$-path in $G$ is a motion plan from configuration $s$ to $t$. \sbmp emerged with the introduction of the \prm algorithm \cite{kslo-prpp-96}. This approach utilizes a roadmap data structure to approximate the connectivity and topology of the free space (\cfree), the set of all valid configurations in \cspace. Motion planning queries are answered by connecting the input start and goal configurations to the roadmap. \rrt \cite{l-rrtnt-98,kl-rrtceasqpp-00} is a single-query algorithm that constructs a search tree in \cspace ~by extending new edges from graph nodes toward random samples. \new{In contrast to real-world physical visibility, \cite{sln-vbprmp-00} utilizes \cspace{} visibility to create sparse roadmaps, rejecting samples in \cspace{} regions already visible by existing nodes.} \sbmp algorithms were found to be effective for a wide range of motion planning problems, and many more variants of \rrt and \prm as well as new \sbmp algorithms were added over the years. See \cite{ock-smpcr-24} for a comparative review of \sbmp algorithms. 
